\section{Mission des Projekts}

\subsection{Erläuterung des zu lösenden Problems}
In jedem Veranstaltungszyklus entstehen hunderte neue Folien,  die ständig angepasst, optimiert und verändet werden. Dort den Überblick zu behalten und effektiv zu arbeiten gestaltet sich häufig als schwierig und erschwert den stressigen Vorlesungsbetrieb weiter.
Die Lösung dafür ist \textbf{Similar Slides}, ein zentrales Repository in dem alle Versionen einer bestimmten .pptx-Datei gesammelt, sortiert und versioniert werden. So kann sichergestellt werden, dass alle Mitarbeitenden einer Veranstaltung immer Zugriff auf die neuste und vergangene Versionen der Folien haben und sich die Unteschiede dieser schnell und klar erkennen lassen.
% Beschreiben Sie in diesem Paragraphen, welche Probleme durch das Projekt gelöst werden sollen. Dies dient als eine erste Einleitung und soll den Lesern ein Grundverständnis von Zweck und Einsatz der Software vermitteln. Einleitungen wie diese schreiben sich am besten, wenn bereits alle anderen Kapitel fertiggestellt wurden, da es sich hierbei um eine leicht vertändliche und kurze Zusammenfassung der folgenden Kapitel handeln soll. Orientieren Sie sich für diesen Abschnitt grob an folgender Struktur:

%   1. Wie heißt das Projekt
%   2. Beschreiben Sie das Problem (Welches Problem? Wann tritt es auf? In welchem Kontext?)
%   3. Beschreiben Sie, wie das Projekt das Problem lösen soll
%   4. Nennen Sie die möglichen Zielgruppen und wie diese von der fertigen Software beeinflusst werden

% Ca. 9-12 Zeilen Text wären hier angebracht.

\subsection{Wünsche und Prioritäten des Kunden}
	Es folgen die Wünsche des Kunden, nach absteigender Priorität geordnet:
	\begin{enumerate}
		\item Es muss eine Repository geben, in das die Foliensätze hochgeladen werden können
        \item Es sollen keine Folien verloren gehen
        \item Foliensätze, die genau so wie gespeichert erneut hochgeladen werden, sollen nicht gespeichert werden
        \item Leicht veränderte Folien sollen nebeneinander gestellt und nach Änderungsdatum sortiert gespeichert werden
        \item Einzelne Folien sollen aus dem Repo in PowerPoint hineinkopiert werden können
        \item Im Repo sollen kleiner Behälter angelegt werden können, in denen Mitarbeitende ihre Folien hochladen können
        \item Die Ähnlichkeit der Folien soll mit einem Ähnlichkeitsmaß gemessen werden
	\end{enumerate}

% nur funktionale anforderungen oder alle?

% Listen Sie hier die vom Kunden genannten Wünsche zum Endprodukt auf und geben Sie an, wie wichtig dem Kunden die verschiedenen Wünsche sind. Eine Sortierung würde sich hier anbieten.

\subsection{Domänenbeschreibung}

Similar Slides ist hauptsächlich auf den Lehrbetrieb einer Hochschule oder Universität ausgelegt. Dort soll es im Veranstaltungs- beziehungsweise Vorlesungsbetrieb unter Microsoft Windows in Kombination mit der PowerPoint Software eingesetzt werden. Dies geschieht während der Ausarbeitung der Folien, also in der Vorbereitung auf die tatsächliche Präsentation. 

% Erläutern Sie, in welcher Domäne das Endprodukt eingesetzt wird. Dazu gehören der Anwendungskontext (wo wird die Software eingesetzt? Z.B. auf dem Hochschulinformationstag) und der technische Kontext (Auf welchen Geräten läuft die Software? Z.B. Mobil, Tablet, Server). Gehen Sie entsprechend darauf ein, welche Effekte die Umgebung auf die Anwendung der Software haben könnte (Lichtverhältnisse, Textgröße, etc.)

% Ca. 8 Zeilen Text wären hier angebracht.
	
\subsection{Maßnahmen zur Anforderungsanalyse}
Zur Erhebungs der Spezifikationen wurde das Video des Kundeninterviews auf dessen Anforderungen untersucht. Anschließend wurden diese in Decision Records festgehalten. Alle offenen Fragen wurden mithilfe des Customer-Proxys geklärt und die Decision Records entsprechend aktualisiert.

% Bitte einfügen: Welche Mockups, Überlegungen, Prototypen wurden erstellt, um die Anforderungen zu erheben? Jeweils nur sehr kurz erläutern, wozu sie dienten, was dabei herausgekommen ist. Sie müssen an dieser Stelle KEINE Mockups etc. erstellen. Überlegen Sie sich, welche Maßnahmen HÄTTEN ergriffen werden können/müssen, um die Anforderungen erfolgreich zu erheben und welchen Zweck diese erfüllt hätten. Hier können auch weniger formale Ansätze gewählt werden, es müssen nicht zwingend die o.g. Mockups, Prototypen sein.